Dado el siguiente problema de programación lineal:

\begin{equation}
\begin{split}
\mbox{max. } z = 2x_1 + 4x_2\\
s.a.:\\
2x_1 + 4x_2 \leq  100\\
1x_1 + 1x_2 \leq 40\\
1x_1 + 3x_2 \leq 80\\
x_1,\,\,x_2,\,\,h_1,\,\,x_2,\,\,h_3\geq 0\\
\end{split}
\end{equation}

Para una solución en la que $x_1 = 10$ y $x_2 = 20$, se pide calcular la inversa de la base correspondiente a la solución dada.

\vspace{1cm}

\noindent
\textbf{Solución}

\vspace{0.5cm}

El problema dado en formato estándar sería:

\begin{equation}
\begin{split}
\mbox{max. } z = 2x_1 + 4x_2\\
s.a.:\\
2x_1 + 4x_2 + h_1 = 100\\
 x_1 +  x_2 + h_2 = 40\\
 x_1 + 3x_2 + h_3 = 80\\
x_1,\,\,x_2\geq 0\\
\end{split}
\end{equation}

Sustituyendo en los valores de $x_1$ y $x_2$ en las restricciones se obtienen los valores de $h_1=0$, $h_2=10$ y $h_3=10$. Es decir, la solución tiene la forma:

\begin{equation}
\begin{split}
x=\begin{pmatrix}
10\\
20\\
0\\
10\\
10\\
\end{pmatrix}
\end{split}
\end{equation}


La solución dada tiene cuatro componentes no nulas y el problema tiene tres restricciones. Las columnas de la matriz $A$ correspondientes a las variables no nulas son linealmente dependientes y, por lo tanto, la solución dada no es básica, por lo que no tiene sentido hablar de base o de inversa de la base.